%%%%%%%%%%%%%%%%%
% This is an sample CV template created using altacv.cls
% (v1.7.4, 30 July 2025) written by LianTze Lim (liantze@gmail.com). Compiles with pdfLaTeX, XeLaTeX and LuaLaTeX.
%
%% It may be distributed and/or modified under the
%% conditions of the LaTeX Project Public License, either version 1.3
%% of this license or (at your option) any later version.
%% The latest version of this license is in
%%    http://www.latex-project.org/lppl.txt
%% and version 1.3 or later is part of all distributions of LaTeX
%% version 2003/12/01 or later.
%%%%%%%%%%%%%%%%

%% Use the "normalphoto" option if you want a normal photo instead of cropped to a circle
% \documentclass[10pt,a4paper,withhyper,normalphoto]{altacv}

\documentclass[10pt,a4paper,withhyper]{altacv}
%% AltaCV uses the fontawesome5 and simpleicons packages.
%% See http://texdoc.net/pkg/fontawesome5 and http://texdoc.net/pkg/simpleicons for full list of symbols.

% Change the page layout if you need to
\geometry{left=1.25cm,right=1.25cm,top=1.5cm,bottom=1.5cm,columnsep=1.2cm}

% The paracol package lets you typeset columns of text in parallel
\usepackage{paracol}

% Change the font if you want to, depending on whether
% you're using pdflatex or xelatex/lualatex
% WHEN COMPILING WITH XELATEX PLEASE USE
% xelatex -shell-escape -output-driver="xdvipdfmx -z 0" sample.tex
\iftutex
  % If using xelatex or lualatex:
  \setmainfont{Roboto Slab}
  \setsansfont{Lato}
  \renewcommand{\familydefault}{\sfdefault}
\else
  % If using pdflatex:
  \usepackage[rm]{roboto}
  \usepackage[defaultsans]{lato}
  % \usepackage{sourcesanspro}
  \renewcommand{\familydefault}{\sfdefault}
\fi

% Change the colours if you want to
\definecolor{SlateGrey}{HTML}{2E2E2E}
\definecolor{LightGrey}{HTML}{666666}
\definecolor{DarkPastelRed}{HTML}{450808}
\definecolor{PastelRed}{HTML}{8F0D0D}
\definecolor{GoldenEarth}{HTML}{E7D192}
\colorlet{name}{black}
\colorlet{tagline}{PastelRed}
\colorlet{heading}{DarkPastelRed}
\colorlet{headingrule}{GoldenEarth}
\colorlet{subheading}{PastelRed}
\colorlet{accent}{PastelRed}
\colorlet{emphasis}{SlateGrey}
\colorlet{body}{LightGrey}

% Change some fonts, if necessary
\renewcommand{\namefont}{\Huge\rmfamily\bfseries}
\renewcommand{\personalinfofont}{\footnotesize}
\renewcommand{\cvsectionfont}{\LARGE\rmfamily\bfseries}
\renewcommand{\cvsubsectionfont}{\large\bfseries}


% Change the bullets for itemize and rating marker
% for \cvskill if you want to
\renewcommand{\cvItemMarker}{{\small\textbullet}}
\renewcommand{\cvRatingMarker}{\faCircle}
% ...and the markers for the date/location for \cvevent
% \renewcommand{\cvDateMarker}{\faCalendar*[regular]}
% \renewcommand{\cvLocationMarker}{\faMapMarker*}


% If your CV/résumé is in a language other than English,
% then you probably want to change these so that when you
% copy-paste from the PDF or run pdftotext, the location
% and date marker icons for \cvevent will paste as correct
% translations. For example Spanish:
% \renewcommand{\locationname}{Ubicación}
% \renewcommand{\datename}{Fecha}


%% Use (and optionally edit if necessary) this .tex if you
%% want to use an author-year reference style like APA(6)
%% for your publication list
% \input{pubs-authoryear}

%% Use (and optionally edit if necessary) this .tex if you
%% want an originally numerical reference style like IEEE
%% for your publication list
\input{pubs-num}

%% sample.bib contains your publications
\addbibresource{sample.bib}

\begin{document}
\name{Ruadhrí Keegan}
\tagline{2nd Year Computer Science Student}
%% You can add multiple photos on the left or right
%\photoR{2.8cm}{Globe_High}
% \photoL{2.5cm}{Yacht_High,Suitcase_High}

\personalinfo{%
  % Not all of these are required!
  \email{roomacaodhagain@gmail.com}
  \phone{+353-086-414-3368}
  \mailaddress{16 Priory Drive, Eden Gate, Delgany}
  \location{Co. Wicklow, Ireland}
  \email{ruadhri.keegan@ucdconnect.ie}
  \linkedin{ruadhri-keegan-2128a9309}
  \github{happy-quinn7}
  %% You can add your own arbitrary detail with
  %% \printinfo{symbol}{detail}[optional hyperlink prefix]
  % \printinfo{\faPaw}{Hey ho!}[https://example.com/]

  %% Or you can declare your own field with
  %% \NewInfoFiled{fieldname}{symbol}[optional hyperlink prefix] and use it
  %%
  %% For services and platforms like Mastodon where there isn't a
  %% straightforward relation between the user ID/nickname and the hyperlink,
  %% you can use \printinfo directly e.g.
  % \printinfo{\faMastodon}{@username@instace}[https://instance.url/@username]
  %% But if you absolutely want to create new dedicated info fields for
  %% such platforms, then use \NewInfoField* with a star:
 % \NewInfoField*{mastodon}{\faMastodon}
  %% then you can use \mastodon, with TWO arguments where the 2nd argument is
  %% the full hyperlink.
  %\mastodon{@username@instance}{https://instance.url/@username}
}

\makecvheader
%% Depending on your tastes, you may want to make fonts of itemize environments slightly smaller
% \AtBeginEnvironment{itemize}{\small}

%% Set the left/right column width ratio to 6:4.
\columnratio{0.6}

% Start a 2-column paracol. Both the left and right columns will automatically
% break across pages if things get too long.
\begin{paracol}{2}
\cvsection{Experience}

\cvevent{Commis Chef}{The Harbour Kitchen}{June 2024 -- Present}{Greystones, Co. Wicklow}
\begin{itemize}
\item Managing multiple sections at once (Starters, Desserts).  
\item Working under pressure during service. Managing responsibilities such as ingredient orders, food preparation, and time management.  
\item Honing a wide range of skills and learning quickly on the job. 
\item Working well in a small team, and working with (and learning from) a wide range of coworkers. 
\item Teaching skills such as food preparation, recipes, and how to work during service to new employees. 
\item Keyholder, responsibility for opening the kitchen.
\end{itemize}

\divider

\cvevent{Additional Experience}{}{}{}
\begin{itemize}
    \item Commis Chef/ Kitchen Porter, The Wicklow Arms (July 2023 - Jan 2024)
    \item Garden Maintenance, Freelance (Jan 2023 - Jun 2023)
    \item Store Assistant, Dunnes Stores (Nov 2022 - Jan 2023)
    \item Store Assistant, Supervalu (Aug 2021 - Jul 2022)
\divider
    \item Continuously employed since age 16 alongside full time education
\end{itemize}

\cvsection{Projects}

\cvevent{Quax - Board Game Implementation}{UCD - Computer Science}{}{}
\begin{itemize}
\item Implemented a fully functional board game in Java, with back-end game logic and error handling, as well as a full GUI.
\item Worked with two of my peers in a Scrum team, worked on this project for a full semester over 4 sprints. The last 2 sprints involved the implementation of an intelligent bot strategy that beats human players approximately 80\% of the time.

\end{itemize}
\divider
\cvevent{Algorithms - Research Paper Analysis}{UCD - Computer Science}{}{}
\begin{itemize}
    \item Implemented a sorting algorithm from a published academic paper in Java, identifying a critical indexing bug absent from the paper's own testing.
    \item Built a benchmarking framework with JVM warm-up handling and 100-trial averaging.
    \item Corrected the proposed algorithm implementation, and when corrected it outperformed Java's QuickSort by 3.6$\times$ at 100,000 elements.
\end{itemize}

\medskip

% use ONLY \newpage if you want to force a page break for
% ONLY the current column

%% Switch to the right column. This will now automatically move to the second
%% page if the content is too long.
\switchcolumn

\cvsection{Strengths}

% Don't overuse these \cvtag boxes — they're just eye-candies and not essential. If something doesn't fit on a single line, it probably works better as part of an itemized list (probably inlined itemized list), or just as a comma-separated list of strengths.

\cvtag{Hard-working}
\cvtag{Calm under Pressure}\\
\cvtag{Analytical Thinker}
\cvtag{Works well in a team}\\
\cvtag{Fast Learner}
\cvtag{Detail-Oriented}

\divider\smallskip

\cvtag{Java}
\cvtag{C}
\cvtag{SQL}
\cvtag{Assembly}
\cvtag{OOP}\\
\cvtag{Data Structures}
\cvtag{Algorithms}\\
\cvtag{Git}
\cvtag{Python}
\cvtag{UNIX}

\cvsection{Languages}

\cvtag{English (Native)}

\cvtag{Irish (Fluent)}

\cvtag{French (Learning)}

 %% Supports X.5 values.

%% Yeah I didn't spend too much time making all the
%% spacing consistent... sorry. Use \smallskip, \medskip,
%% \bigskip, \vspace etc to make adjustments.
\smallskip

\cvsection{Education}

\cvevent{BSc Computer Science}{University College Dublin}{Sept 2024 -- Present}{}
Current GPA: 3.7/4.2

\divider

\cvevent{Leaving Certificate}{Temple Carrig School}{Sept 2022 -- June 2024}{}
\begin{itemize}
    \item Irish - H1 (90-100)
    \item Maths - H1
    \item English - H2 (80-89)
    \item Computer Science - H1
    \item Politics and Society - H1
    \item Physics - H2
    \item Design and Comms. Graphics - H3 (70-79)
\end{itemize}
% \divider
\cvsection{Referees}

% \cvref{name}{email}{mailing address}
\cvref{Prof.\ Barry Smyth}{UCD - Director of BSc Computer Science}{bscdirector.cs@ucd.ie}{}
\divider
\cvref{Chef Jaco Pretorius}{The Harbour Kitchen}{+353-096-788-6559}
{}

\end{paracol}

\end{document}
